% Full title: Main Results and Proof Architecture
\section{Main Results and Proof Architecture}
\label{sec:main-results}

We collect the assumptions required by the target theorem.  They define the
\emph{regular parameter set} $\mathfrak R$ of pairs $(k,\beta)$ for which:
\begin{enumerate}[label=(R\arabic*),leftmargin=2.8em]
  \item $k>0$ is real and $k^2$ lies in the compact gap window
  $I_\beta$ of \eqref{eq:strict-gap};
  \item $k\notin\mathfrak E_{\mathrm{rep}}(\beta)$ and all layer, Rayleigh,
  and port-trace maps in \Cref{sec:muller-coupling} are bounded on the stated
  spaces;
  \item the half-guide translations have no unit-circle spectrum and
  \Cref{prop:relation-Bloch} holds with injective Riesz synthesis maps
  $\mathcal E^\pm$;
  \item the center representation in \Cref{thm:center-representation} is
  surjective and $\mathcal N_{\mathrm c}(k,\beta)$ is closed.
\end{enumerate}
These assumptions deliberately exclude thresholds, Wood anomalies, embedded
guided states, complex resonances, and representation poles.

Before stating the main equation, it is useful to keep its three objects at
different conceptual levels:
\begin{enumerate}[label=(\alph*),leftmargin=2.2em]
  \item $\mathbb A(k,\beta):\mathcal X(k,\beta)\to\mathcal Y(k,\beta)$ is the
  \emph{residual operator} defined in \eqref{eq:coupled-operator}.  Its input
  is one tuple $x=(\eta,\xi,c^-,c^+)$ of center interface densities,
  homogeneous Rayleigh coefficients, and left/right stable-mode
  coefficients.  The equation $\mathbb A x=0$ says exactly that the center
  transmission jumps and both center--half-guide Cauchy mismatches vanish.
  \item $\mathcal N_{\mathrm{rep}}(k,\beta)$ is the subset of
  $\ker\mathbb A(k,\beta)$ defined in \eqref{eq:Nrep-definition} whose
  densities and modal coefficients reconstruct the zero field everywhere.
  It is removed because two algebraic descriptions that generate the same
  physical field must represent the same solution class.
  \item $\mathcal G(k,\beta)=\ker(A(\beta)-k^2I)$ is the physical guided-mode
  eigenspace defined in \eqref{eq:guided-space}: its elements are global,
  vertically $\beta$-quasiperiodic $H^1$ solutions on the infinite strip and
  are therefore square-integrable in the transverse direction.
\end{enumerate}
Thus the theorem does not identify all algebraic vectors with fields one by
one; it identifies algebraic kernel vectors modulo their zero-field
representations with the physical eigenspace.

\begin{theorem}[Kernel--guided-mode equivalence: main target]
\label{thm:main-kernel-equivalence}
Let $(k,\beta)\in\mathfrak R$.  The reconstruction
$\mathcal R_{\mathrm g}(k,\beta)$ maps
$\ker\mathbb A(k,\beta)$ onto the global guided-mode eigenspace
$\mathcal G(k,\beta)$, and
\begin{equation}
  \ker\!\left(
  \mathcal R_{\mathrm g}(k,\beta)|_{\ker\mathbb A(k,\beta)}
  \right)
  =\mathcal N_{\mathrm{rep}}(k,\beta)
  =\mathcal N_{\mathrm c}(k,\beta)\times\{0\}\times\{0\}.
  \label{eq:representation-kernel-identification}
\end{equation}
Consequently, reconstruction induces a natural linear isomorphism
\begin{equation}
  \overline{\mathcal R}_{\mathrm g}(k,\beta):
  \frac{\ker\mathbb A(k,\beta)}{\mathcal N_{\mathrm{rep}}(k,\beta)}
  \xrightarrow{\ \cong\ }\mathcal G(k,\beta).
  \label{eq:main-quotient-isomorphism}
\end{equation}
In particular,
\begin{equation}
  \dim\!\left(
  \ker\mathbb A(k,\beta)/\mathcal N_{\mathrm{rep}}(k,\beta)
  \right)=m_{\mathrm g}(k,\beta).
  \label{eq:multiplicity-preservation}
\end{equation}
If the stronger density-injectivity statement
$\mathcal N_{\mathrm c}(k,\beta)=\{0\}$ holds, then the quotient can be
removed and $\ker\mathbb A(k,\beta)\cong\mathcal G(k,\beta)$.
\end{theorem}
\status{main proof target.}
\begin{proofroadmap}
For completeness, take $u\in\mathcal G(k,\beta)$.  By
\Cref{thm:bounded-reduction}, its center restriction belongs to
$\mathcal S_{\mathrm c}(k,\beta)$ and its half-guide restrictions lie in the two
outgoing spaces.  Apply \Cref{thm:center-representation} to choose a center
class $[\eta,\xi]$, and apply \eqref{eq:relation-stable-range} to obtain the
unique generalized-mode coordinates $c^\pm$.  The physical interface and port
matching conditions then give $x=(\eta,\xi,c^-,c^+)\in\ker\mathbb A$.

For soundness, start with $x\in\ker\mathbb A$.  The first row of
\eqref{eq:coupled-operator} makes the center reconstruction a physical
transmission field.  The other rows identify its complete Cauchy data with
those of the two synthesized half-guide fields.  Apply
\Cref{lem:weak-gluing} and the strict-gap $H^1$ selection to obtain a member of
$\mathcal G(k,\beta)$.

For injectivity, assume the reconstructed global field is zero.  Injectivity
of the generalized Riesz synthesis gives $c^-=c^+=0$.  The center pair then
lies in $\mathcal N_{\mathrm c}$, proving
\eqref{eq:representation-kernel-identification}.  This last step is the main
proof bottleneck if the half-guide synthesis is only a closed-range map or if the
center density kernel has not been characterized.  The closest prior
ingredients are \cite[Theorem~4.5]{Fliss2013},
\cite[Theorem~4 and Appendix~A]{BarnettGreengard2010}, and
\cite[Theorem~2.2]{HohageSoussi2013}; none by itself proves the coupled
statement.

The full internal obligation is headed ``continuous kernel--guided-space
isomorphism'' in
\path{research/planning/muller-cauchy/p-proofs.md}; the dependency
analysis is in
\path{research/planning/muller-cauchy/p-deps.md}.  The
weakest acceptable conclusion is a surjection from the quotient kernel onto
$\mathcal G(k,\beta)$ together with an abstract, rather than explicit,
description of $\mathcal N_{\mathrm{rep}}$.
\end{proofroadmap}
\begin{proof}
Fix $(k,\beta)\in\mathfrak R$ and suppress these parameters from the
notation.

\emph{Step 1: reconstruction maps the algebraic kernel into the guided-mode
space.}
Let $x=(\eta,\xi,c^-,c^+)\in\ker\mathbb A$.  Its first block equation is
$\mathcal M(\eta,\xi)=0$.  Hence the two material traces of
$\mathcal R_{\mathrm c}(\eta,\xi)$ agree, and the center reconstruction is a
field in $\mathcal S_{\mathrm c}$.  The two remaining block equations say
\[
  \Pi^\pm(\eta,\xi)=Q^\pm c^\pm
  =\Tr^{0,\pm}\mathcal E^\pm c^\pm.
\]
Thus the center field and each synthesized half-guide field have identical
center-oriented Cauchy data at their common port.  Since
$\mathcal E^\pm c^\pm\in\mathcal U^\pm_{\mathrm{out}}$, two applications of
\Cref{lem:weak-gluing} show that $\mathcal R_{\mathrm g}x$ is a global field
in $H^1_\beta(B)$ satisfying the weak eigenvalue equation.  By
\Cref{prop:weak-strong},
\[
  \mathcal R_{\mathrm g}x\in\mathcal G(k,\beta).
\]
Consequently the restriction of $\mathcal R_{\mathrm g}$ to
$\ker\mathbb A$ is a well-defined linear map into $\mathcal G(k,\beta)$.

\emph{Step 2: every guided mode has algebraic coordinates.}
Take $u\in\mathcal G(k,\beta)$.  By
\Cref{thm:bounded-reduction}, its center restriction
$u^0:=u|_{\mathcal C^0}$ is a physical center solution and the restrictions
$u^\pm:=u|_{W^\pm}$ belong to
$\mathcal U^\pm_{\mathrm{out}}$.  Assumption (R4) and
\Cref{thm:center-representation} provide
$(\eta,\xi)\in\mathcal Z_{\mathrm c}$ such that
\[
  \mathcal R_{\mathrm c}(\eta,\xi)=u^0.
\]
By assumption (R3) and \eqref{eq:relation-stable-range}, there are
$c^\pm\in\mathcal K_s^\pm$ such that
\[
  Q^\pm c^\pm=\Tr^{0,\pm}u^\pm.
\]
The converse part of \Cref{lem:weak-gluing} gives
$\Tr^{0,\pm}u^0=\Tr^{0,\pm}u^\pm$.  Therefore
$x=(\eta,\xi,c^-,c^+)$ makes all three rows of
\eqref{eq:coupled-operator} vanish.  It lies in $\ker\mathbb A$, and its
piecewise reconstruction is exactly $u$.  Hence
$\mathcal R_{\mathrm g}|_{\ker\mathbb A}$ is onto
$\mathcal G(k,\beta)$.

\emph{Step 3: identify the representation kernel.}
The first equality in \eqref{eq:representation-kernel-identification} follows
immediately from the definition \eqref{eq:Nrep-definition}.  We prove the
second.  Suppose $x\in\ker\mathbb A$ and
$\mathcal R_{\mathrm g}x=0$.  On the two half-guides this says
\[
  \mathcal E^-c^-=0,
  \qquad
  \mathcal E^+c^+=0.
\]
The injectivity required in (R3) gives $c^-=c^+=0$.  On the center cell,
$\mathcal R_{\mathrm c}(\eta,\xi)=0$, while the first row of
$\mathbb A x=0$ gives $(\eta,\xi)\in\mathcal Z_{\mathrm c}$.  Thus
$(\eta,\xi)\in\mathcal N_{\mathrm c}$ and
\[
  \mathcal N_{\mathrm{rep}}
  \subset\mathcal N_{\mathrm c}\times\{0\}\times\{0\}.
\]
Conversely, if $(\eta,\xi)\in\mathcal N_{\mathrm c}$ and $c^-=c^+=0$,
then $\mathcal M(\eta,\xi)=0$ and
$\Pi^\pm(\eta,\xi)=\Tr^{0,\pm}\mathcal R_{\mathrm c}(\eta,\xi)=0$.
Hence $(\eta,\xi,0,0)\in\ker\mathbb A$ and reconstructs the zero global
field.  The reverse inclusion follows, proving
\eqref{eq:representation-kernel-identification}.

\emph{Step 4: quotient and multiplicity.}
We have proved that the linear map
$\mathcal R_{\mathrm g}|_{\ker\mathbb A}$ is surjective and has kernel
$\mathcal N_{\mathrm{rep}}$.  The first isomorphism theorem gives
\eqref{eq:main-quotient-isomorphism}.  Taking dimensions and using
\eqref{eq:geometric-multiplicity} yields
\eqref{eq:multiplicity-preservation}.  Finally, if
$\mathcal N_{\mathrm c}=\{0\}$, then
\eqref{eq:representation-kernel-identification} makes the restricted
reconstruction injective, so no quotient is needed.
\end{proof}

\subsection{Interpretation and the Fredholm question}

The term ``spurious-free'' is justified at the continuous level only after
\eqref{eq:representation-kernel-identification}: a kernel class then carries a
nonzero guided field unless it is the zero class.  This does not control large
inverse norms or discretization-induced near-kernels.  Those distinct issues
are emphasized by the spurious quasi-resonance results in
\cite[Theorems~1.10--1.12]{HiptmairMoiolaSpence2022}.

If $\mathcal N_{\mathrm c}$ is closed, let
\[
  \widehat{\mathcal X}(k,\beta):=
  \mathcal X(k,\beta)/
  \bigl(\mathcal N_{\mathrm c}(k,\beta)\times\{0\}\times\{0\}\bigr).
\]
Since this nullspace is contained in $\ker\mathbb A$, the coupled operator
descends to a bounded map
$\widehat{\mathbb A}(k,\beta):\widehat{\mathcal X}(k,\beta)
\to\mathcal Y(k,\beta)$.  Recall that a bounded operator between Banach spaces
is Fredholm if its range is closed and its kernel and cokernel are finite
dimensional; its index is $\ind F=\dim\ker F-\dim\operatorname{coker}F$.

\begin{conjecture}[Fredholm index-zero realization]
\label{conj:Fredholm}
There is a natural square multitrace or quotient realization of
$\widehat{\mathbb A}(k,\beta)$ on the regular parameter set for which
\begin{equation}
  \widehat{\mathbb A}(k,\beta)
  =\mathbb A_0+\mathbb K(k,\beta),
  \qquad \mathbb A_0^{-1}\text{ bounded},\quad
  \mathbb K(k,\beta)\text{ compact},
  \label{eq:Fredholm-splitting}
\end{equation}
and hence $\widehat{\mathbb A}(k,\beta)$ is Fredholm of index zero.
\end{conjecture}
\status{currently conjectural strengthening.}
\begin{proofroadmap}
The principal candidate for $\mathbb A_0$ combines the M\"uller/Calder\'on
principal block with identity blocks obtained by Riesz-normalizing the stable
trace syntheses.  Smooth quasiperiodic-minus-free-space kernel differences and
trace maps between separated curves are compact candidates.  Every remainder
must be checked on the precise Sobolev products in
\eqref{eq:coupled-domain}--\eqref{eq:coupled-codomain}; a block display that is
not square does not have a Fredholm index.

See \path{research/planning/muller-cauchy/notes/fredholm.md} and
the obligations named ``reference operator plus compact perturbation'' and
``Fredholm index zero'' in
\path{research/planning/muller-cauchy/p-proofs.md}.  The closest
structural precedents are \cite[Theorem~4]{BarnettGreengard2010} and
\cite[Lemma~5.1]{HohageSoussi2013}.  If no natural square realization exists,
the conjecture must be replaced by a closed-range or semi-Fredholm statement,
or by an index-zero theorem only after finite modal truncation.  The main
quotient equivalence in \Cref{thm:main-kernel-equivalence} does not logically
depend on this conjecture.
\end{proofroadmap}
